\documentclass[lettersize,journal]{IEEEtran}
\usepackage{amsmath,amsfonts,amssymb}
\usepackage{amsthm}
\usepackage{algorithm}
\usepackage{algorithmic}
\usepackage{array}
\usepackage[caption=false,font=normalsize,labelfont=sf,textfont=sf]{subfig}
\usepackage{textcomp}
\usepackage{stfloats}
\usepackage{url}
\usepackage{graphicx}
\usepackage{booktabs,multirow}
\usepackage[numbers]{natbib}
\usepackage{xcolor}
\usepackage[colorlinks,allcolors=blue]{hyperref}
\newtheorem{proposition}{Proposition}
\hyphenation{op-tical net-works semi-conduc-tor IEEE-Xplore}

\begin{document}

\title{One Policy Does Not Fit All: Matching Patient Streaming Policies to Emergency Department Types}

\author{Farhood~Rismanchian, Farzad~Zaerpour, and Peyman~Hosseinzadeh~Kassani% author list and affiliations to confirm
\thanks{Manuscript received XX; revised XX.}}

\markboth{IEEE Transactions on Engineering Management}%
{Rismanchian \MakeLowercase{\textit{et al.}}: One Policy Does Not Fit All}

\maketitle

\begin{abstract}
To be written.
\end{abstract}

\begin{IEEEkeywords}
Emergency department, patient streaming, prioritization, discrete-event simulation, policy fit, healthcare operations management.
\end{IEEEkeywords}

% Section plan follows notes/PROJECT_BRIEF.md
\section{Introduction}
\IEEEPARstart{E}{mergency} departments differ in the patients they serve and in how they are organized.
% Motivation: one-size-fits-all streaming policies; research questions; contributions.

\section{Literature Review}
% ED streaming and prioritization rules; APQ; prediction-based dispatching; policy comparisons across ED types; ED performance indicators.

\section{Simulation Framework}
\subsection{Patient Generator and ED Types}
% Synthetic patients; true vs predicted attributes; pathway sampling based on \citet{Rismanchian2023_PatientJourney}.
\subsection{Process Model and Resources}
\subsection{Streaming Policies}
% Queue disciplines and bed-assignment rules; similarity-based representation.
\subsection{Policy Fidelity of the Similarity Representation}
% Propositions + decision-level and outcome-level fidelity tests.

\section{Experimental Design}
% Factors and levels; KPIs; warm-up, run length, replications, CRN; verification and validation (STRESS-DES).

\section{Results}
% Best policy by ED type; regret of one-size-fits-all; trade-offs; value of prediction accuracy.

\section{Managerial Implications}

\section{Limitations and Future Research}

\section{Conclusion}

\bibliographystyle{IEEEtranN}
\bibliography{ed_library}

\end{document}
